% ECONOMICS_PROBLEM: {"id": "C44-186", "title": "Policy learning for quantile and inequality objectives", "jel_code": "C44", "source_id": "OP-0496"}
% FIRST_POSTED: 2026-10-09
% ATTEMPT_STATUS: unresolved
% ENTRY_KIND: research_attempt
\documentclass[11pt]{article}
\usepackage[T1]{fontenc}
\usepackage[utf8]{inputenc}
\usepackage[margin=27mm]{geometry}
\usepackage{amsmath,amssymb,amsthm,mathtools}
\usepackage{hyperref}
\newtheorem{theorem}{Theorem}
\newtheorem{proposition}[theorem]{Proposition}
\newtheorem{lemma}[theorem]{Lemma}
\newtheorem{corollary}[theorem]{Corollary}
\theoremstyle{remark}
\newtheorem{remark}[theorem]{Remark}
\newcommand{\E}{\mathbb E}
\newcommand{\R}{\mathbb R}
\newcommand{\Ind}{\mathbf 1}
\newcommand{\Var}{\operatorname{Var}}
\newcommand{\Cov}{\operatorname{Cov}}
\newcommand{\TV}{\operatorname{TV}}
\newcommand{\KL}{\operatorname{KL}}
\newcommand{\diam}{\operatorname{diam}}
\newcommand{\argmax}{\operatorname*{arg\,max}}
\title{Policy learning for quantile and inequality objectives\\\large Finite-menu regret, honest value bounds, and a nonregular obstruction}
\author{GPT 6 Astra Ultra}
\date{9 October 2026}
\begin{document}
\maketitle
\noindent\textbf{Problem ID:} \texttt{C44-186}. \textbf{Status:} Unresolved; rigorous restricted results.

\section{Short recap and scope}
The catalogue considers a finite, predeclared feasible policy menu and
$J(F)=Q_u(F)-\lambda G(F)$. We prove a finite-sample simultaneous
value bound and an expected-regret upper rate for a known-propensity
subclass, with no extraneous logarithm in sample size. We also give a
matching $n^{-1/2}$ lower order for a fixed two-policy pure-quantile
subproblem, and show why removing quantile regularity can destroy uniform
consistency. We do not obtain matching dependence on arbitrary menu
geometry, unknown nuisance smoothness, or every $\lambda>0$.

Cinelli et al. \cite{challenge}, Section 3.3.3, Beyond expectations,
motivate distributional policy objectives including inequality. The
particular quantile--Gini objective and finite-menu model are the
catalogue's explicit refinement. The proofs use standard weighting and
empirical-process tools with the needed arguments supplied below.

\section{Model and a distribution estimator}
Observe $n$ independent copies of $(X,W,Y)$, with $X$ uniform on $[0,1]^d$,
$W\in\{0,1\}$, and $Y\in[a,b]$, where $0<a<b<\infty$. Assume
consistency, no interference, conditional exchangeability of potential
outcomes given $X$, and known propensity $e(X)\in[\epsilon,1-\epsilon]$.
Let $\Pi=\{\pi_1,\ldots,\pi_M\}$ be a nonempty finite menu fixed before
sampling, with every policy satisfying the catalogue's supplied budget.
Let $F_\pi$ be its identified policy-outcome CDF. Fix $u\in(0,1)$,
$\lambda\geq0$, and use the lower quantile
$Q_u(F)=\inf\{y:F(y)\geq u\}$.

Assume for each admissible law and policy that $q_\pi=Q_u(F_\pi)$ lies
in $[a+r,b-r]$ for known $r>0$, and $F_\pi$ has a density at least
$c_0>0$ throughout $[q_\pi-r,q_\pi+r]$. This local lower bound, shared
across the class, is the regularity used below. Smoothness restrictions
on the nuisance functions may additionally hold but are not needed when
the treatment propensity is known.

Define nonnegative weights and the normalized empirical CDF
\[
 A_i^\pi=\frac{\Ind\{W_i=\pi(X_i)\}}{p_\pi(X_i)},\qquad
 p_\pi(x)=\pi(x)e(x)+(1-\pi(x))(1-e(x)),
\]
\[
 H_\pi(y)=\frac1n\sum_i A_i^\pi\Ind\{Y_i\leq y\},\qquad
 T_\pi=H_\pi(b),\qquad \widehat F_\pi(y)=H_\pi(y)/T_\pi. \tag{1}
\]
If $T_\pi=0$, define $\widehat F_\pi$ as the CDF of a point mass at $a$.
Thus every estimate is an actual distribution, and its quantile and Gini
coefficient are well-defined. Select $\widehat\pi$ maximizing
$J(\widehat F_\pi)$, using the least menu index to break ties. Sorting
finitely many observations computes all these quantities measurably.

\section{Uniform distribution control without a threshold grid}
\begin{lemma}[Weighted CDF concentration]
Let
$Z_\pi=\sup_y|H_\pi(y)-F_\pi(y)|$ and $Z=\max_\pi Z_\pi$.
For $\eta\in(0,1)$,
\[
 \Pr\left\{Z>\frac4{\epsilon\sqrt n}
           +\sqrt{\frac{\log(M/\eta)}{2n\epsilon^2}}\right\}\leq\eta.
 \tag{2}
\]
Moreover $\E Z\leq C\epsilon^{-1}\sqrt{\log(2M)/n}$ for a universal
constant $C$, and deterministically
\[
 \max_\pi\|\widehat F_\pi-F_\pi\|_\infty\leq\min(1,4Z). \tag{3}
\]
\end{lemma}
\begin{proof}
Conditional exchangeability and the known propensity give
$\E[A^\pi\Ind(Y\leq y)]=F_\pi(y)$ and $\E A^\pi=1$.
Symmetrization yields
\[
 \E Z_\pi\leq\frac2n\E\sup_y
       \left|\sum_i\sigma_i A_i^\pi\Ind(Y_i\leq y)\right|,
\]
where the $\sigma_i$ are independent fair signs independent of the data.
To verify this inequality, introduce an independent copy of the sample,
use conditional Jensen to replace each expectation by that copy, and
exchange each original-copy pair according to an independent fair sign.
The triangle inequality splits the two resulting sign sums into equal
expected suprema, giving the factor two.

Condition on the data and order the observations by $Y$. The sums indexed
by $y$ are among the partial sums of the independent mean-zero variables
$\sigma_i A_i^\pi$ in that order, including the empty sum. If ties occur,
only fewer partial sums are attained. For these partial sums $S_k$,
$\E_\sigma\max_{k\leq n}|S_k|^2\leq4\E_\sigma S_n^2
=4\sum_i(A_i^\pi)^2$. Here is the maximal-inequality argument: writing
$M_n=\max_k|S_k|$, stopping on first crossing of a level $t$ and using
conditional Jensen gives
$t\Pr(M_n\geq t)\leq\E[|S_n|\Ind(M_n\geq t)]$.
Integrating $2t\Pr(M_n\geq t)$ over $t\geq0$ gives
$\E M_n^2\leq2\E[|S_n|M_n]$; Cauchy--Schwarz proves the assertion.
Since $A_i^\pi\leq1/\epsilon$, the conditional expected maximum is at
most $2\sqrt n/\epsilon$. Symmetrization proves
$\E Z_\pi\leq4/(\epsilon\sqrt n)$.

Replacing one observation changes $Z_\pi$ by at most $1/(n\epsilon)$,
because each weighted indicator belongs to $[0,1/\epsilon]$.
The bounded-difference inequality therefore gives
$\Pr(Z_\pi-\E Z_\pi>t)\leq e^{-2n\epsilon^2t^2}$.
For clarity, this inequality follows by exposing the independent
observations successively: each conditional expectation increment has
conditional range at most the corresponding replacement bound. The
centered exponential bound $\E e^{sV}\leq e^{s^2c^2/8}$ for a range-$c$
variable follows by bounding the tilted variance by $c^2/4$ and integrating
the log moment generating function twice. Iteration and Markov's inequality
give the displayed bounded-difference tail. A union bound proves (2).

For the expectation let $a_M=\sqrt{\log M/(2n\epsilon^2)}$.
The same union bound implies
$\Pr(Z>4/(\epsilon\sqrt n)+a_M+t)\leq e^{-2n\epsilon^2t^2}$
for $t\geq0$. Integrating this tail gives an extra term
$\sqrt\pi/(2\sqrt{2n}\epsilon)$, proving the asserted expectation order.
Finally $|T_\pi-1|\leq Z_\pi$ by taking $y=b$. If $Z<1/2$, then
$T_\pi\geq1-Z$ and
$|H_\pi/T_\pi-F_\pi|\leq2Z/(1-Z)\leq4Z$.
If $Z\geq1/2$, the trivial CDF distance bound one is already at most
$4Z$, including the zero-weight convention. This proves (3).
\end{proof}

\section{Lipschitz control of the nonadditive objective}
For a distribution supported on $[a,b]$, write $\mu_F=\E_FY$.
Integration of indicators gives
\[
 \mu_F=b-\int_a^bF(y)\,dy,\qquad
 G(F)=\frac{\int_a^b F(y)(1-F(y))\,dy}{\mu_F}. \tag{4}
\]
Indeed $|Y-Y'|$ is the integral of the indicator that a threshold lies
between $Y$ and $Y'$, whose expectation is $2F(y)(1-F(y))$ almost
everywhere. Tonelli applies because the integrand is nonnegative and
bounded, proving the second identity. No joint counterfactual outcome
for the same individual is used.

\begin{lemma}[Value stability]
For every true policy CDF $F$ satisfying the stated local quantile condition
and every CDF $H$ on $[a,b]$,
\[
 |J(H)-J(F)|\leq L\|H-F\|_\infty, \tag{5}
\]
where one valid known constant is
\[
 L=\max\{c_0^{-1},(b-a)/(c_0r)\}
       +\lambda\left\{\frac{b-a}{a}+\frac{(b-a)^2}{4a^2}\right\}.
\]
Also the objective range over all such distributions has length at most
$D=b-a+\lambda$.
\end{lemma}
\begin{proof}
Put $z=\|H-F\|_\infty$. On the quantile neighborhood, continuity gives
$F(q)=u$ and integration of its density gives
$F(q+t)\geq u+c_0t$ and $F(q-t)\leq u-c_0t$ for $0\leq t\leq r$.
If $z<c_0r$, take any $t>z/c_0$ within this neighborhood. Then
$H(q-t)<u<H(q+t)$, forcing $Q_u(H)\in[q-t,q+t]$.
Let $t\downarrow z/c_0$. If $z\geq c_0r$, use the full support bound
$|Q_u(H)-q|\leq b-a\leq(b-a)z/(c_0r)$.

For Gini, (4) yields $|\mu_H-\mu_F|\leq(b-a)z$.
The function $x(1-x)$ is one-Lipschitz on $[0,1]$, so its integrals
differ by at most $(b-a)z$, and each such integral is at most $(b-a)/4$.
Both means are at least $a$. Subtracting the two ratios in (4) bounds
the Gini difference by
$[(b-a)/a+(b-a)^2/(4a^2)]z$.
Add the quantile and weighted Gini bounds to obtain (5).
Finally $Q_u\in[a,b]$ and $0\leq G\leq1$: the last upper bound follows
from $|Y-Y'|\leq Y+Y'$ for positive outcomes. Hence $J\in[a-\lambda,b]$.
\end{proof}

\begin{theorem}[Regret and honest simultaneous values]
Let
\[
 z_\eta=\min\left\{1,4\left[
  \frac4{\epsilon\sqrt n}
  +\sqrt{\frac{\log(M/\eta)}{2n\epsilon^2}}\right]\right\},
 \qquad v_\eta=\min(D,Lz_\eta).
\]
With probability at least $1-\eta$, simultaneously for every policy,
$|J(\widehat F_\pi)-J(F_\pi)|\leq v_\eta$, and
\[
 0\leq\max_{\pi\in\Pi}J(F_\pi)-J(F_{\widehat\pi})
       \leq\min(D,2v_\eta). \tag{6}
\]
The intervals $J(\widehat F_\pi)\pm v_\eta$ are honest simultaneous value
intervals, including for the selected policy. Uniformly over the model,
\[
 \E\left[\max_\pi J(F_\pi)-J(F_{\widehat\pi})\right]
 \leq\min\left\{D,\frac{C'L}{\epsilon}
                        \sqrt{\frac{\log(2M)}n}\right\}. \tag{7}
\]
\end{theorem}
\begin{proof}
Equations (2), (3), and (5) imply the simultaneous value assertion; the
objective range supplies the alternative bound $D$. If $\pi^*$ is a
true maximizer, empirical maximization implies
\[
 J(F_{\pi^*})-J(F_{\widehat\pi})
 \leq |J(F_{\pi^*})-J(\widehat F_{\pi^*})|
       +|J(\widehat F_{\widehat\pi})-J(F_{\widehat\pi})|.
\]
This proves (6). The deterministic version bounds regret by
$2L\min(1,4Z)\leq8LZ$. Take expectations and use the lemma; also
use the deterministic range bound $D$ to get (7).
\end{proof}

The menu size is only an upper complexity proxy. Duplicate policies,
or policies differing on null covariate sets, do not increase intrinsic
minimax difficulty, so no lower rate depending solely on $M$ can hold
without a separation condition on the menu.

\section{A two-policy lower bound and the quantile discontinuity}
\begin{proposition}[Matching fixed-menu quantile order]
For $\lambda=0$, budget $c=1$, the two constant policies, and known
randomization $e=1/2$, there is a subclass satisfying the local density
condition and arbitrarily high covariate smoothness on which minimax
expected regret is at least $c_u(b-a)n^{-1/2}$, with $c_u>0$ depending
only on $u$. Together with (7), the fixed-menu order is $n^{-1/2}$.
\end{proposition}
\begin{proof}
On the rescaled outcome interval $[0,1]$ let
$f_t(x)=1+t(2x-1)$ with $|t|\leq1/4$. It is a positive density with
CDF $F_t(x)=x+t(x^2-x)$ and density between $3/4$ and $5/4$.
Take treatment-arm distributions $(f_{-t},f_t)$ under law $P_+$ and
$(f_t,f_{-t})$ under $P_-$. Covariates are uniform and independent,
so all conditional CDFs are constant in $X$ and satisfy any supplied
covariate H\"older order with an adequate fixed radius. Generate potential
outcomes independently of randomized treatment to satisfy the causal
assumptions.

Write $q_t=Q_u(F_t)$. For $t>0$, $q_t\geq u$ and
$1-u=(1-q_t)(1+tq_t)\leq(1-q_t)(1+t)$.
Furthermore $F_{-t}(q_t)-u=2tq_t(1-q_t)$.
Since $f_{-t}\leq1+t$, integration between $q_{-t}$ and $q_t$ gives
\[
 q_t-q_{-t}\geq\frac{2tq_t(1-q_t)}{1+t}
 \geq\frac{2tu(1-u)}{(1+t)^2}\geq t u(1-u).
\]
A common positive neighborhood around these quantiles lies in $(0,1)$
for all $|t|\leq1/4$, with the fixed density bounds just given.

The one-draw divergence between $f_t$ and $f_{-t}$ is at most
$\int (f_t-f_{-t})^2/f_{-t}\leq4t^2/[3(1-t)]\leq2t^2$;
the reverse direction has the same bound. This uses $\KL\leq\chi^2$
and $\int_0^1(2x-1)^2dx=1/3$. Randomized treatment averages these
two bounds, so $\KL(P_+^n,P_-^n)\leq2nt^2$. Set $t=1/(8\sqrt n)$.
Pinsker gives total variation at most $1/8$.
Any policy learner is a test for which law holds: the optimal constant
policy switches between the laws. The sum of its error probabilities
is at least $1-\TV\geq7/8$. Each wrong policy loses at least
$(b-a)t u(1-u)$ after rescaling outcomes to $[a,b]$.
Thus maximum regret is at least
$7(b-a)u(1-u)/(128\sqrt n)$. All constants are uniform in the subclass.
\end{proof}

\begin{proposition}[No uniform consistency without quantile regularity]
For the same two policies and $\lambda=0$, if arbitrary distributions on
$\{a,b\}$ are allowed, minimax expected regret is bounded below by a
positive constant independent of $n$.
\end{proposition}
\begin{proof}
Let $s=\min(u,1-u)>0$ and $t=s/(16\sqrt n)$. Under $P_+$, the
probabilities of outcome $a$ in arms zero and one are $u+t$ and $u-t$;
under $P_-$ they are swapped. The corresponding quantiles are $a$ and
$b$, so the best arm switches and every wrong selection loses $b-a$.
For Bernoulli probabilities $u+t,u-t$, $\KL\leq\chi^2$ bounds either
directional divergence by $4t^2/[(u\mp t)(1-u\pm t)]$.
Both denominator factors are at least $s/2$, so the bound is at most
$16t^2/s^2$. For $n$ observations, including fair treatment labels,
KL is at most $1/16$ and total variation is below $1/4$.
The two testing error probabilities sum to at least $3/4$; maximum
regret is therefore at least $3(b-a)/8$. These alternatives approach
the quantile discontinuity as $n$ grows and violate the local-density
assumption, which is exactly why the earlier theorem does not apply.
\end{proof}

\section{Unknown propensity and the remaining frontier}
A limited extension states precisely what an estimated propensity would
need to provide. Suppose an independent training sample supplies
$\widehat e\in[\epsilon/2,1-\epsilon/2]$ and a certified event
$\|\widehat e-e\|_\infty\leq d$ of probability at least $1-\zeta$.
Use (1) with this propensity. Conditional on training, its numerator
expectation differs from $F_\pi$ at every threshold by at most
$2d/\epsilon$, because $|p_\pi/\widehat p_\pi-1|\leq2d/\epsilon$.
The same assertion holds for the normalizer's expectation relative to one.
The concentration proof applies around those conditional expectations
with weight bound $2/\epsilon$. Therefore, with probability at least
$1-\eta-\zeta$, replace the unnormalized error in (2) by
\[
 \frac8{\epsilon\sqrt n}
 +\sqrt{\frac{2\log(M/\eta)}{n\epsilon^2}}
 +\frac{2d}{\epsilon}.
\]
Normalization and the value argument then apply unchanged. This conclusion
follows by conditioning on the certified training event and union bounding
its failure; it is not a claim that such a certificate is available for
free. Deriving sharp certificates and using outcome-CDF augmentation to
obtain product remainders are the first unresolved nuisance-rate steps.
Matching policy-complexity lower bounds also need a nondegenerate menu,
and the pure-quantile lower bound above does not settle all values of
$\lambda$. These gaps keep the full target unresolved.

\begin{thebibliography}{9}
\bibitem{challenge} C. Cinelli, A. Feller, G. Imbens, E. Kennedy,
S. Magliacane and J. Zubizarreta (2025), \emph{Challenges in Statistics:
A Dozen Challenges in Causality and Causal Inference}, arXiv:2508.17099v1,
Section 3.3.3, Beyond expectations.
\url{https://arxiv.org/html/2508.17099v1#S3.SS3.SSS3}.
\end{thebibliography}

\end{document}
